\section{The local determinant}\label{sec:local-determinant}

\subsection{An integral for the determinant}\label{subsec:determinant-integral}

The determinant in \eqref{eq:Gaussian-block-integration} has a local
logarithm obtained by inserting the polymer propagators into a
resolvent integral. We first derive the integral at real connections.
Put $n_k=\ell^3$ and
\[
 c_k^d=4\bigl((1-L^{-3})\log b_k+L^{-3}\log(b_k+b_L)\bigr)
       \quad(k\ge1),\qquad c_0^d=0.
\]
For real $A\in\widetilde{\mathcal G}_k$ and $0<\mu\le1$,
\begin{align}
 \delta Z_k(A)
  &=\exp\left\{c_k^dn_k-ib_k^2\int_0^\infty\dd y\,
       \operatorname{Tr}\bigl[\gamma_5\mathcal B_{k,y}
          \mathcal Q_kS_y\mathcal Q_k^\sharp\mathcal B_{k,y}\bigr]
                        \right\}\quad(k\ge1),
                                    \label{eq:determinant-integral-positive-level}\\
 \delta Z_0(A)
  &=\exp\left\{\int_1^\infty\frac{\dd y}{y}
                       \operatorname{Tr}[F_0S_y]
             -i\int_0^1\dd y\,\operatorname{Tr}[\gamma_5S_y]
                        \right\}.\label{eq:determinant-integral-level-zero}
\end{align}
The trace in the first line is on the unit spinor space. At level
zero it is the normalized chamber trace
$\operatorname{Tr}O=48^{-1}\sum_{c,g}\tr O((c,g),(c,g))$.

To prove these identities, write $K=K_k$ as in the preceding
subsection. At real $A$, the fine adjoint identity
$D_A^*=\gamma_5D_A\gamma_5$, the reality of $\mu,b_k,b_L$, and
the orthogonal averaging projections give
$K^*=\gamma_5K\gamma_5$. Hence $K\gamma_5$ is Hermitian and
invertible. Since $\det\gamma_5=1$ on each four-component spinor,
its determinant is $\det K$. Let $\lambda_j$ be its real nonzero
eigenvalues, counted with multiplicity. There are $4n_k$ at positive
levels and $192n_0$ at level zero. Then
\[
 \operatorname{Tr}\bigl[\gamma_5(K+i\gamma_5y)^{-1}\bigr]
        =\sum_j(\lambda_j+iy)^{-1},\qquad
 \operatorname{Tr}\bigl[K(K+i\gamma_5y)^{-1}\bigr]
        =\sum_j\lambda_j(\lambda_j+iy)^{-1}.
\]
For a real $\lambda\ne0$, choose along $y\ge0$ the logarithm
with value $\log|\lambda|+i\pi\mathbf1[\lambda<0]$ at zero.
Its derivative is $i/(\lambda+iy)$ and
$\operatorname{Log}(\lambda+iy)-\log y\to i\pi/2$.
Integration gives
\begin{align*}
 -i\int_0^Y\frac{\dd y}{\lambda+iy}
        &=\operatorname{Log}\lambda-
                          \operatorname{Log}(\lambda+iY),\\
 \int_1^\infty\frac{\dd y}{y}\frac{\lambda}{\lambda+iy}
       -i\int_0^1\frac{\dd y}{\lambda+iy}
        &=\operatorname{Log}\lambda-i\pi/2.
\end{align*}
Summing the second identity proves
\eqref{eq:determinant-integral-level-zero}: its constant phase is
$-96\pi i n_0\in2\pi i\Z$.

For $k\ge1$, the ranks of $P_L$ and $1-P_L$ are respectively
$4L^{-3}n_k$ and $4(1-L^{-3})n_k$. The explicit formula for
$\mathcal B_{k,y}$ therefore gives
\[
 -i\int_0^Y\dd y\,\operatorname{Tr}[\gamma_5\mathcal B_{k,y}]
          =-4n_k\log Y+c_k^dn_k+o(1)
                  \quad(Y\to\infty).
\]
Insert \eqref{eq:coarse-resolvent-formula} in the first scalar
integral. The terms $4n_k\log Y$ cancel. The remaining phase is
$2\pi i n_k$, which disappears on exponentiation, and this proves
\eqref{eq:determinant-integral-positive-level}.
The integral converges absolutely: block compression bounds
$\mathcal Q_kS_y\mathcal Q_k^\sharp$ uniformly, while each row
and column sum of $\mathcal B_{k,y}$ is at most
$(1+2\exp(1))(b_k^2+y^2)^{-1/2}$. Its trace is bounded by
\[
 128\exp(1)(1+2\exp(1))^2a'n_k(b_k^2+y^2)^{-1}.
\]
At level zero the spectral identities already give absolute
convergence. The uniform kernel bounds below will give a bound
independent of the eigenvalues.

\subsection{Local determinant terms}\label{subsec:local-determinant-terms}

Choose $\gamma_{\mathrm w}\le\gamma$ and
$\kappa_{\mathrm w}=\kappa+\kappa_0+3$ with $\kappa,\kappa_0>0$.
The propagator bounds above give a fixed constant $C_w$ such that,
in component sup norms, for $y\ge0$,
\begin{align*}
 \|\mathbf1_{\Delta_z}S_{k,y}(X,A)\mathbf1_{\Delta_{z'}}\|
       &\le C_we^{-\gamma_{\mathrm w}d(z,z')
                                  -\kappa_{\mathrm w}d_M(X)}
                                                    \quad(k\ge1),\\
 |S_{0,y}(X,A)(x,x')|_c
       &\le C_w(1+y)^{-1}e^{-\gamma_{\mathrm w}d(c,c')
                                  -\kappa_{\mathrm w}d_M(X)}.
\end{align*}
Here $|\cdot|_c$ is the matrix norm induced by the component sup
norm. We may use the same $C_w$ after reducing the rates from
$\gamma,\kappa'$.
For $X\in\mathcal P_k$ define
\begin{align}
 \widetilde E_k(X,A)
    &=-ib_k^2\int_0^\infty\dd y\,
       \operatorname{Tr}\bigl[\gamma_5\mathcal B_{k,y}\mathcal Q_k
           S_{k,y}(X,A)\mathcal Q_k^\sharp\mathcal B_{k,y}\bigr]
                   \quad(k\ge1),\label{eq:local-determinant-definition}\\
 \widetilde E_0(X,A)
    &=\int_1^\infty\frac{\dd y}{y}
                     \operatorname{Tr}[F_0S_{0,y}(X,A)]
       -i\int_0^1\dd y\,\operatorname{Tr}[\gamma_5S_{0,y}(X,A)],\notag\\
 E_k^d(X,A)&=\tfrac14\bigl(\widetilde E_k(X,A)
       +\widetilde E_k(X,-A)
       +\overline{\widetilde E_k(X,\overline A)}
       +\overline{\widetilde E_k(X,-\overline A)}\bigr),\notag\\
 E_k^{\mathrm{det}}(X,A)&=E_k^d(X,A)-E_k^d(X,0).\notag
\end{align}
Set $c_X=54$ and
\[
 C_d=16c_XC_w\max\left\{\tfrac\pi2b\exp(1)(1+2\exp(1)),\,
                                  2+58\exp(1)+b\exp(1)\right\}.
\]

\begin{proposition}\label{prop:local-determinant-expansion}
The integrals \eqref{eq:local-determinant-definition} converge
absolutely for $A\in\widetilde{\mathcal G}_k$. They define analytic
functions, depending only on bonds meeting $X$, with
\begin{equation}\label{eq:local-determinant-bound}
 |\widetilde E_k(X,A)|, |E_k^d(X,A)|
               \le C_dM^3e^{-(\kappa+1)d_M(X)}.
\end{equation}
They satisfy
\begin{equation}\label{eq:local-determinant-quotient}
 \delta Z_k(A)=e^{c_k^dn_k+\sum_X\widetilde E_k(X,A)},\qquad
 \frac{\delta Z_k(A)}{\delta Z_k(0)}
                         =e^{\sum_XE_k^{\mathrm{det}}(X,A)}.
\end{equation}
The functions $E_k^d$ are gauge invariant, even in $A$, and satisfy
$\overline{E_k^d(X,\overline A)}=E_k^d(X,A)$. They are covariant
under the lattice step group.
\end{proposition}

\begin{proof}
For $k\ge1$ write
$N_y(X)=\mathcal Q_kS_{k,y}(X,A)\mathcal Q_k^\sharp$.
It is supported on the unit sites of $X$, and each entry is bounded
by $\exp(1)C_we^{-\kappa_{\mathrm w}d_M(X)}$ by normalized block
compression. To estimate the trace, cycle one factor
$\mathcal B_{k,y}$ to the left and use
\[
 G_y=\mathcal B_{k,y}\gamma_5\mathcal B_{k,y}
     =\gamma_5\bigl((b_k+i\gamma_5y)^{-2}(1-P_L)
              +(b_k+b_L+i\gamma_5y)^{-2}P_L\bigr).
\]
Its column sum in Euclidean matrix norm is at most
$(1+2\exp(1))/(b_k^2+y^2)$: the identity contributes one and the two
projection coefficients contribute at most $2\exp(1)$. Since
$|\tr(mm')|\le8|m|_2\max_{i,j}|m'_{ij}|$ for four-component
matrices, the integrand is bounded by
\begin{equation}\label{eq:polymer-determinant-trace-bound}
 8\exp(1)(1+2\exp(1))C_wM^3|X|_M
           e^{-\kappa_{\mathrm w}d_M(X)}(b_k^2+y^2)^{-1}.
\end{equation}
Integrating and multiplying by $b_k^2$ leaves the factor
$\pi b_k/2\le\pi b/2$.

At $k=0$, the formula for $D_A$ and the averaging penalty give
the scalar component row bound
\begin{equation}\label{eq:first-step-Dirac-row-bound}
 \sup_{x,\beta}\sum_{\beta'}\int\dd x'\,
                    |F_0(A)(x,x')_{\beta\beta'}|
                      \le R_F:=58\exp(1)+1+b\exp(1).
\end{equation}
The diagonal and three internal hops contribute
at most $41+\mu$, the external hop at most
$6e^{\mathfrak i(A)}$, and the block penalty at most $b\exp(1)$.
The bound follows from $\mu\le1$ and
$e^{\mathfrak i(A)}\le \exp(1)$. Point evaluation has normalized
kernel $48\mathbf1[x=x']$; its factor $48$ cancels the chamber
measure. The two traces in the level-zero integral are consequently
bounded by
\[
 4R_FC_wM^3|X|_M(1+y)^{-1}e^{-\kappa_{\mathrm w}d_M(X)},
 \qquad
 4C_wM^3|X|_M(1+y)^{-1}e^{-\kappa_{\mathrm w}d_M(X)}.
\]
The second bound uses that $\gamma_5$ has four entries of modulus
one. Their integrals contribute respectively
$\int_1^\infty(y(1+y))^{-1}\dd y=\log2$ and
$\int_0^1(1+y)^{-1}\dd y=\log2$.

The geometric bound \eqref{eq:polymer-size-tree} gives
$|X|_M\le c_X(d_M(X)+1)\le c_Xe^{d_M(X)}$.
Applying this to the trace bounds,
and using $\kappa_{\mathrm w}-1\ge\kappa+1$, proves
\eqref{eq:local-determinant-bound} for $\widetilde E_k$ and hence
for its four-term mean. The majorants are independent of $A$ and
integrable in $y$. Each integrand is analytic in $A$; Cauchy's
formula on a compact polydisk and the majorant permit its contour
integrals to be exchanged with the $y$ integral. The result is
analytic.

Every phase in the compressed trace belongs to a unit block or an
$L$-block meeting $X$, and that block is contained in $X$ by
nesting. The fine propagator is supported on $X$ and uses only its
connection. At level zero a nonzero external-hop trace uses a bond
meeting $X$; the penalty uses the $L$-block of its row. This proves
locality. Gauge transformations conjugate each product under the
trace and hence leave it unchanged. A lattice transformation whose
spin matrix commutes with $\gamma_5$ does the same. Every spatial
action of the step group has such a lift: multiplication by
$\gamma_1\gamma_2\gamma_3$ changes its commutation sign with
$\gamma_5$ and leaves its spatial action fixed. This proves step
covariance. The mean in \eqref{eq:local-determinant-definition}
then makes $E_k^d$ even and real in the stated analytic sense.

For real connections, sum the polymer propagators and use
\eqref{eq:determinant-integral-positive-level} or
\eqref{eq:determinant-integral-level-zero}. This proves the first
identity of \eqref{eq:local-determinant-quotient} there. Both sides
are analytic on $\widetilde{\mathcal G}_k$, which is open and
star-shaped. They agree on a real neighbourhood of zero, so their
Taylor series agree there and analytic continuation proves the
identity throughout the domain.

To justify the symmetrized quotient, put
$\mathcal L(A)=\sum_X\widetilde E_k(X,A)$.
The charge-conjugation transpose and the complex-conjugation
identity give
$\delta Z_k(-A)=\delta Z_k(A)
=\overline{\delta Z_k(\overline A)}$.
For the latter, the spin matrix $\varkappa=\gamma_5\mathsf C$
conjugates $\overline{K_k(\overline A)}$ to $K_k(-A)$.
Thus $\mathcal L(-A)-\mathcal L(A)$ and
$\overline{\mathcal L(\overline A)}-\mathcal L(A)$ take values
in $2\pi i\Z$. They are continuous on the connected domain,
so they are constant. Subtracting their values at zero gives
\[
 \sum_X\bigl(E_k^d(X,A)-E_k^d(X,0)\bigr)
                            =\mathcal L(A)-\mathcal L(0).
\]
This proves the second identity in
\eqref{eq:local-determinant-quotient} without a choice of logarithm
branch for the complex determinant.
\end{proof}

\subsection{The small-field gain for the determinant}\label{subsec:determinant-small-field-gain}

For a photon connection write $A_\mu(c)=A(c,c+\eta_ke_\mu)$ and
\[
 \partial_\nu A_\mu(c)
       =\eta_k^{-1}(A_\mu(c+\eta_ke_\nu)-A_\mu(c)),\qquad
 \|\delta_\alpha\partial A\|_\infty
       =\sup_{\mu,\nu}\sup_{0<d(c,c')\le1}
          \frac{|\partial_\nu A_\mu(c')-\partial_\nu A_\mu(c)|}
               {d(c,c')^\alpha}.
\]
Local versions use only differences and bonds inside the given set.
In this section the curl is mesh normalized:
$(dA)_{\mu\nu}=\partial_\mu A_\nu-\partial_\nu A_\mu$.
Let $R_k$ be the complex bond functions whose three norms
$\|A\|_\infty$, $\|\partial A\|_\infty$, and
$\|\delta_\alpha\partial A\|_\infty$ are less than
$e_k^{-3/4+\varepsilon}$, $e_k^{-3/4+2\varepsilon}$, and
$e_k^{-3/4+3\varepsilon}$, respectively. Let $\widetilde R_k$
consist of $A_0+A_1$, with $A_1\in R_k$ and $A_0$ real, such that
on every $\square^\dagger=\square^{\sim L}$ one can write
$A_0=A'+d\lambda$, where $\lambda,A'$ are real and the same
three local bounds hold for $A'$. The external connection domain is
\[
 \mathfrak X_k=\{\mathscr C:\|\mathscr C\|_\infty
                  <e_k^{-1+2\varepsilon},\quad
             \|\partial\mathscr C\|_\infty<e_k^{-1+3\varepsilon}\}.
\]
If $e_k^{1/4}+e_k^\varepsilon<1$, both $\mathfrak X_k$ and
$\widetilde R_k+\mathfrak X_k$ lie in
$\widetilde{\mathcal G}_k$: on the local cubes the real part has
a representative of size $e_k^{-3/4+\varepsilon}$, and the
remaining part has size at most
$(e_k^{1/4}+e_k^\varepsilon)e_k^{-1+\varepsilon}$.

Put $v_k=-\log e_k$. Assume
\begin{equation}\label{eq:determinant-gain-conditions}
 v_k\ge9,\qquad e_k^\varepsilon\le\tfrac14,\qquad
 8M(v_k/4+4)e_k^\varepsilon\le1,\qquad
 \ell\ge Mv_k^2.
\end{equation}

\begin{proposition}\label{prop:determinant-small-field-gain}
For $\mathscr A\in\widetilde R_k$, $\mathscr C\in\mathfrak X_k$,
and $X\in\mathcal P_k$, assume \eqref{eq:determinant-gain-conditions}
and $e_k^{1/4}+e_k^\varepsilon<1$. Then
\begin{equation}\label{eq:determinant-small-field-gain}
 |E_k^d(X,\mathscr A+\mathscr C)-E_k^d(X,\mathscr C)|
       \le8C_dM^3e_k^{1/4}e^{-\kappa d_M(X)}.
\end{equation}
\end{proposition}

\begin{proof}
If $d_M(X)>v_k/4$, bound the two terms separately by
\eqref{eq:local-determinant-bound} and use
$e^{-d_M(X)}\le e_k^{1/4}$. Suppose therefore that
$d_M(X)\le v_k/4$. The polymer lift and diameter estimates give
a nonwrapping lattice box $W$ containing every bond meeting $X$,
with each side of length at most
$s_W=M(d_M(X)+3)$.
Their applicability follows from
$2M(d_M(X)+2)<\ell$ and $M(d_M(X)+4)<\ell$, which hold by
\eqref{eq:determinant-gain-conditions}.

Write $\mathscr A=A_0+A_1$ as in the domain definition.
On each plaquette, the local real representative gives
$\|dA_0\|_\infty<2e_k^{-3/4+2\varepsilon}$.
Choose a fixed-order rectilinear path $\Gamma(o,c)$ from a site
$o\in W$ and set $\lambda(c)=\eta_k A_0(\Gamma(o,c))$ on $W$.
Then
\[
 \|A_0-d\lambda\|_{\infty,W}
             \le2s_W\|dA_0\|_{\infty,W}
             <4s_We_k^{-3/4+2\varepsilon}.
\]
To see the factor $2s_W$, append one bond to $\Gamma(o,c)$ and
subtract $\Gamma(o,c')$. After cancelling common bonds, the
resulting loop bounds at most two rectangles of width $\eta_k$.
Together they contain at most $2s_W/\eta_k$ plaquettes, each
contributing $\eta_k dA_0$. The paths and rectangles remain in
$W$ because it is a product of intervals. Extend the real function
$\lambda$ by zero outside the projected box. On every bond meeting
$X$, the bound is unchanged. By
\eqref{eq:determinant-gain-conditions} it is at most
$\tfrac12e_k^{-3/4+\varepsilon}$.

Let $B=A_0-d\lambda+A_1$ on bonds meeting $X$ and zero elsewhere.
Then $\|B\|_\infty<\tfrac32e_k^{-3/4+\varepsilon}$.
For $|t|\le\tfrac13e_k^{-1/4}$,
\[
 \|tB+\mathscr C\|_\infty
       <(\tfrac12+e_k^\varepsilon)e_k^{-1+\varepsilon}
       \le\tfrac34e_k^{-1+\varepsilon}.
\]
Thus this connection belongs to the complex connection ball
$\mathcal G_k=\{A:\|A\|_\infty<e_k^{-1+\varepsilon}\}$,
which is contained in $\widetilde{\mathcal G}_k$, and
$f(t)=E_k^d(X,tB+\mathscr C)$ is analytic on this disk and
bounded by the right side of \eqref{eq:local-determinant-bound}.
Gauge invariance and locality identify the difference in
\eqref{eq:determinant-small-field-gain} with $f(1)-f(0)$.
On the circle of radius $r=\tfrac14e_k^{-1/4}>2$, Cauchy's
formula yields
\[
 f(1)-f(0)=\frac1{2\pi i}\oint_{|t|=r}
                         \frac{f(t)}{t(t-1)}\,\dd t,
 \qquad |f(1)-f(0)|\le\frac{\sup|f|}{r-1}
                          \le8e_k^{1/4}\sup|f|.
\]
This proves the claimed gain, retaining the stronger decay
$e^{-(\kappa+1)d_M(X)}$ in this case.
\end{proof}

\begin{proof}[Proof of \Cref{thm:fermion-block-integration}]
\Cref{prop:fine-coarse-resolvents} and
\eqref{eq:level-zero-covariance} give invertibility on the two
integration spaces. The weighted completion calculation gives
\eqref{eq:Gaussian-block-integration}, with the first-step weight from
\eqref{eq:first-step-Grassmann-weight}.
\Cref{prop:local-determinant-expansion} gives the analytic determinant
terms, their bound, symmetries and normalized quotient.
\end{proof}


